\section{Revisiting the Problem Statement}
\label{sec:disc-problem-statement}

The problem statement asked how to build a researcher-operated adaptive reading platform in which sensing, eye-movement analysis, decision strategy, and intervention execution can each be replaced independently of the others, while the platform still runs real Tobii reading sessions and applies context-preserving micro-interventions without disturbing the participant's reading (\cref{sec:final-problem-statement}).
We answer this one research question at a time.
For each question we judge how far the goal was reached and state the caveat that limits it.
We give the verdicts here and leave the wider design knowledge behind them to \cref{sec:disc-design-reflections}.
Throughout, we point to the evidence in \cref{ch:evaluation} instead of repeating its figures.

\paragraph{RQ1 (Modularity): met strongly.}
The first research question asked whether sensing, analysis, decision, and intervention can be separated behind stable contracts, so that a new module can be added without changing the core reading runtime.
The evidence says yes, and \cref{sec:eval-modularity} shows it at three levels: the build enforces the direction of dependencies, adding a module changes only one place, and external providers connect through the documented contract.
The strongest case is the analysis pipeline of a concurrent, independent master's thesis at DTU, connected to our analysis port through the documented contract (\cref{sec:eval-modularity}).
It integrates a real research system we did not write, against a contract we did not design around it, and adding it changed no file in the backend.
More than connect, this external provider was the live analysis source for the whole evaluation: the oculomotor events we report were produced by an outside team's algorithm rather than by ours.
This is the closest evidence we have that the contract works for someone other than its authors, and that it works well enough for an independent team to run their own study through it.
The build-enforced boundary is what sets this result apart from the earlier project prototypes, which were extended by copying the previous frontend instead of building on a shared contract (\cref{subsec:related-rtr-prior-work}).
Here, a reference that breaks the boundary fails to compile, so modularity is a rule the build checks rather than a habit the authors must keep.
The caveat that remains, developed in \cref{sec:disc-limitations}, is narrower than it first appears: our \emph{reference} decision and analysis providers were built by us and so show only that the contract is sufficient, but the independent thesis integration, being the analysis source of the evaluation rather than a one-off connection test, is the evidence that the seam holds beyond our own work.

\paragraph{RQ2 (Sensing pipeline): met for sensing and transport.}
The second research question asked whether a real Tobii stream can be turned into oculomotor events fast enough for live adaptation, in a way that is inspectable and reproducible.
For the sensing and transport stages, the answer is yes.
The tracker streamed at its rated rate with high gaze validity, every recorded sample stayed within the latency budget set in \cref{sec:non-functional-requirements}, and the stages of the pipeline can be inspected and replayed from the exported record (\cref{sec:eval-performance}).
A narrower gap remains on the automated decision path.
Its latency is instrumented both inside the backend and out to an external provider, and the eight behavioural sessions, run in manual mode, exercised neither leg.
A dedicated advisory-mode validation run did exercise both, with a reference provider proposing changes that the console gated, and both the in-process pipeline and the out-of-process round trip stayed within budget (\cref{sec:eval-performance}).
The budget is therefore now shown for the automated decision path as well, though on a single short run with a rule-based provider on one host rather than across a campaign with a real decision model.
We see the remaining gap as one of dataset scale and model realism, not of architecture, and describe the campaign that would close it in \cref{sec:disc-future-work}.

\paragraph{RQ3 (Context-preserving intervention): the architecture is met, and the restore strategy is a measured finding.}
The third research question asked how a typographic micro-intervention can change the text while keeping the reader's place, the recovery concern raised in the literature \cite{jensen2025context}.
The evidence has two parts that must be kept separate.
The first part is the architecture.
Its claim was that the cost of an intervention to the reader's continuity should be measured and recorded, not assumed, and this claim is met.
Every intervention wrote a graded outcome and its residual error into the session record, and on this small, descriptive sample the two behavioural proxies, reading-resume time and post-intervention regression, both moved in the expected direction, lower with preservation than without (\cref{sec:eval-context}).
The second part is the specific strategy, and here the same record shows a weakness: on small reflows the semantic-restart strategy moved the reader too far, further than the reflow itself had moved them.
That the reader nonetheless resumed faster despite this larger displacement is best explained by the highlight cue that accompanies the restore, which flashes the restored sentence and draws the eye straight back to the resume point, so the behavioural gain is plausibly carried by the highlight rather than by the geometry of the restore, an account consistent with the prior project finding that highlighting was the component that drove the significant recovery improvement \cite{jensen2025context}.
This reading reconciles the two measures rather than opposing them: the highlight restores the reader's place perceptually even where the positional arithmetic overshoots, which sharpens rather than softens the case for improving the restore geometry, and it is the confound recorded in \cref{sec:eval-threats}.
We read these findings together as support for the design choice, not as a problem with it.
The architecture was built to make this cost visible, and it reported an unfavourable result about one strategy instead of hiding it, which is exactly what the measurable-disruption idea of \cref{sec:contributions} promised.
Why the over-repositioning happens, and how a future restore strategy could avoid it, is covered in \cref{sec:disc-lessons} and \cref{sec:disc-future-work}.

\paragraph{RQ4 (Researcher control and experimentability): met.}
The fourth research question asked whether the platform gives the researcher the control and the auditable record, including session replay, needed to run experiments and to compare interventions and decision strategies.
The recorded sessions were run from start to finish through a single console that is gated on setup and calibration.
The whole of \cref{ch:evaluation} was then rebuilt from the exported records alone, with no access to any live state, which shows the auditable record in practice rather than only claiming it (\cref{sec:eval-control}).
Session replay rebuilds a recorded session in the reading interface from its exported record alone (FR21), and this is the reproducibility that the project's planned comparison workflow depends on.
This is the gap the earlier prototypes left open.
Each of them was a self-contained build and could not support systematic comparison, whereas the auditable, replayable record makes comparing interventions and strategies a feature of the platform rather than separate work for every study (\cref{subsec:gap-experimentability}).

Together, the four answers close the loop back to the three gaps that motivated the thesis and to the contribution claimed at the start.
The experimentability gap (\cref{subsec:gap-experimentability}) is closed by the gated console and the auditable, replayable record (RQ4).
The sensing-coupling gap (\cref{subsec:gap-sensing}) is closed by the sensing adapter that lets the eye tracker and the mouse source be swapped for each other (RQ1 and RQ2).
The external-decision integration gap (\cref{subsec:gap-integration}) is closed by the provider contract, which lets external and AI-driven deciders connect as a built-in property of the architecture rather than as code added to this repository (RQ1).
On this basis we judge the main contribution of \cref{sec:contributions} supported: one system brings together real-time, context-preserving typographic intervention, pluggable decision logic, researcher-in-the-loop control, and reproducible session capture, a combination that the landscape of \cref{sec:market-landscape} leaves empty.
The honest qualifier is that the decision logic is pluggable and was used through manual operation and reference providers, but its automated path was not run end to end on real hardware in this work, which the limitations and future work discuss.
